\begin{lemma}[协变导数的交换子]
对光滑规范势及光滑截面，规范协变导数的交换子为曲率在截面上的作用，带电场的具体系数按\eqnrefs{eq:covariant-derivative}固定，结果为\eqnrefs{eq:covariant-commutator}。
\end{lemma}

\begin{proof}
先把协变性要求写成对任意测试旋量都成立的算符等式
\[
D'_\mu\,U_\chi=U_\chi D_\mu,
\qquad
U_\chi(x)=\ee^{\ii q_{\mathrm e}\chi(x)/\hbar}.
\]
若局域联络仍取一阶形式
\[
D'_\mu=\partial_\mu+\frac{\ii q_{\mathrm e}}{\hbar}A'_\mu,
\]
则作用在任意 $\psi$ 上给
\begin{align*}
D'_\mu(U_\chi\psi)
={}&U_\chi\partial_\mu\psi
+\frac{\ii q_{\mathrm e}}{\hbar}U_\chi(\partial_\mu\chi)\psi
+\frac{\ii q_{\mathrm e}}{\hbar}A'_\mu U_\chi\psi,\\
U_\chi D_\mu\psi
={}&U_\chi\partial_\mu\psi
+\frac{\ii q_{\mathrm e}}{\hbar}U_\chi A_\mu\psi.
\end{align*}
由于这两式必须对任意 $\psi$ 相等，乘法算符的系数必须逐点相等：
\[
A'_\mu+\partial_\mu\chi=A_\mu,
\]
即正文\eqnrefs{eq:gauge-A-lower}。因此在给定电荷约定和一阶局域导数结构后，$A_\mu$ 的变换律不是额外选择，而是协变性唯一固定的。

再计算曲率时，不能只把 $D_\mu$ 当成普通数相乘；应让对易子作用在任意测试场上：
\begin{align*}
[D_\mu,D_\nu]\psi
={}&\frac{\ii q_{\mathrm e}}{\hbar}
\bigl[(\partial_\mu A_\nu)-(\partial_\nu A_\mu)\bigr]\psi\\
&+\left(\frac{\ii q_{\mathrm e}}{\hbar}\right)^2
(A_\mu A_\nu-A_\nu A_\mu)\psi.
\end{align*}
所有 $A\,\partial\psi$ 交叉项成对消失。对阿贝尔 $U(1)$，$A_\mu$ 是普通乘法函数，所以第二行也消失，留下
\[
[D_\mu,D_\nu]\psi
=\frac{\ii q_{\mathrm e}}{\hbar}F_{\mu\nu}\psi,
\qquad
F_{\mu\nu}=\partial_\mu A_\nu-\partial_\nu A_\mu.
\]
最后代入 $A'_\mu=A_\mu-\partial_\mu\chi$，有
\[
F'_{\mu\nu}-F_{\mu\nu}
=-\partial_\mu\partial_\nu\chi
+\partial_\nu\partial_\mu\chi=0,
\]
其中只用了光滑标量函数的混合偏导可交换。这同时解释了为什么阿贝尔场强是联络的一阶反对称导数，以及为什么非阿贝尔推广中会多出联络自身的对易子项。
\end{proof}

\begin{lemma}[动力学动量的交换子]
在本章 SI 制最小耦合约定下，$\pi_i=-\ii\hbar\partial_i-qA_i$ 满足 $[\pi_i,\pi_j]=\ii q\hbar\epsilon_{ijk}B_k$，其中 $\vb{B}=\nabla\times\vb{A}$。
\end{lemma}

\begin{proof}
正文\eqnrefs{eq:kinetic-momentum} 已写出动力学动量与正则动量、矢势之间的关系。不同空间分量的非对易性因此完全来自 $[\hat p_i,A_j]$；把位置表象的基本对易关系逐项展开，再用 $\vb{B}=\vb{\nabla}\times\vb{A}$ 收集反对称导数，就得到正文\eqnrefs{eq:pi-commutator}。


由 $\hat\Pi_i=\hat p_i-q_{\mathrm e} A_i(\vb{r},t)$ 与 $[\hat p_i,f(\vb{r},t)]=-\ii\hbar\partial_i f$，逐项展开
\begin{align*}
 [\hat\Pi_i,\hat\Pi_j]
 &=[\hat p_i-q_{\mathrm e} A_i,\hat p_j-q_{\mathrm e} A_j]\\
 &=[\hat p_i,\hat p_j]
 -q_{\mathrm e}[\hat p_i,A_j]
 -q_{\mathrm e}[A_i,\hat p_j]
 +q_{\mathrm e}^2[A_i,A_j].
\end{align*}
坐标表象中的普通阿贝尔势是乘法算符，因此 $[A_i,A_j]=0$，而正则动量满足 $[\hat p_i,\hat p_j]=0$。剩余两项为
\begin{align*}
 [\hat\Pi_i,\hat\Pi_j]
 &=-q_{\mathrm e}(-\ii\hbar\partial_iA_j)
 -q_{\mathrm e}(+\ii\hbar\partial_jA_i)\\
 &=\ii q_{\mathrm e}\hbar(\partial_iA_j-\partial_jA_i).
\end{align*}
利用 $B_k=\epsilon_{k\ell m}\partial_\ell A_m$ 以及
\begin{align*}
 \epsilon_{ijk}B_k
 &=\epsilon_{ijk}\epsilon_{k\ell m}\partial_\ell A_m\\
 &=(\delta_{i\ell}\delta_{jm}-\delta_{im}\delta_{j\ell})
 \partial_\ell A_m\\
 &=\partial_iA_j-\partial_jA_i,
\end{align*}
得到主文公式。
\end{proof}

\begin{proposition}[最小耦合 Dirac 算子的平方]
在光滑电磁背景下，利用 Clifford 关系平方 Dirac 算子，得到二阶协变波动算子及由 $\sigma^{\mu\nu}F_{\mu\nu}$ 给出的自旋项，完整系数见\eqnrefs{eq:dirac-squared}。此式为微分表达式恒等式，闭算子定义域另行指定。
\end{proposition}

\begin{proof}
正文\eqnrefs{eq:dirac-em} 是含任意时空依赖电磁场的一阶 Dirac 方程。左乘符号相反的一阶因子后，协变导数乘积的对称部分给出 $D_\mu D^\mu$，反对称部分则通过 $[D_\mu,D_\nu]$ 化成场强并与 $\sigma^{\mu\nu}$ 收缩，从而得到正文\eqnrefs{eq:dirac-squared}。


将 \eqnrefs{eq:dirac-em} 除以 $c$，
\begin{equation*}
 (\ii\hbar\slashed D-m_{\mathrm e} c)\psi=0.
\end{equation*}
左乘共轭的一阶因子 $\ii\hbar\slashed D+m_{\mathrm e} c$：
\begin{align*}
 0
 &=(\ii\hbar\slashed D+m_{\mathrm e} c)(\ii\hbar\slashed D-m_{\mathrm e} c)\psi\\
 &=\left[(\ii\hbar)^2\gamma^\mu\gamma^\nu D_\mu D_\nu-m_{\mathrm e}^2c^2\right]\psi\\
 &=\left[-\hbar^2\gamma^\mu\gamma^\nu D_\mu D_\nu-m_{\mathrm e}^2c^2\right]\psi.
\end{align*}
由 $\sigma^{\mu\nu}=\ii[\gamma^\mu,\gamma^\nu]/2$ 与 Clifford 代数，
\begin{align*}
 \gamma^\mu\gamma^\nu
 &=\frac12\{\gamma^\mu,\gamma^\nu\}
 +\frac12[\gamma^\mu,\gamma^\nu]\\
 &=\eta^{\mu\nu}\id()-\ii\sigma^{\mu\nu}.
\end{align*}
因此
\begin{align*}
 \gamma^\mu\gamma^\nu D_\mu D_\nu
 &=D_\mu D^\mu-\ii\sigma^{\mu\nu}D_\mu D_\nu.
\end{align*}
由于 $\sigma^{\mu\nu}$ 对 $\mu,\nu$ 反对称，$D_\mu D_\nu$ 的对称部分与它收缩为零：
\begin{align*}
 \sigma^{\mu\nu}D_\mu D_\nu
 &=\frac12\sigma^{\mu\nu}
 (D_\mu D_\nu-D_\nu D_\mu)\\
 &=\frac12\sigma^{\mu\nu}[D_\mu,D_\nu]\\
 &=\frac{\ii q_{\mathrm e}}{2\hbar}\sigma^{\mu\nu}F_{\mu\nu}.
\end{align*}
于是
\begin{align*}
 \gamma^\mu\gamma^\nu D_\mu D_\nu
 &=D_\mu D^\mu
 +\frac{q_{\mathrm e}}{2\hbar}\sigma^{\mu\nu}F_{\mu\nu}.
\end{align*}
代回二阶方程并整体乘以 $-1$：
\begin{align*}
 0
 &=\left[
 \hbar^2D_\mu D^\mu+m_{\mathrm e}^2c^2
 +\frac{q_{\mathrm e}\hbar}{2}\sigma^{\mu\nu}F_{\mu\nu}
 \right]\psi,
\end{align*}
即\eqnrefs{eq:dirac-squared}。
\end{proof}

\begin{proposition}[偶奇算符的二阶块对角化]
设 $H=\beta mc^2+\mathcal E+\mathcal O$，$[\beta,\mathcal E]=0$、$\{\beta,\mathcal O\}=0$。在本节静态背景及低能弱混合展开条件下，逐阶消去奇部得到\eqnrefs{eq:fw-even-operator-form-r16}；其适用阶数由所保留的嵌套交换子决定。
\end{proposition}

\begin{proof}
取第一步反厄米生成元

\begin{equation*}
 A_1\equiv\ii S_1=\frac{\beta_D\hat{\mathcal O}}{2m_{\mathrm e} c^2},
 \qquad U_1=\ee^{A_1}.
\end{equation*}
由于 $\{\beta_D,\hat{\mathcal O}\}=0$，有 $(\beta_D\hat{\mathcal O})^\dagger=-\beta_D\hat{\mathcal O}$，所以 $U_1$ 幺正。Baker--Campbell--Hausdorff 展开为
\begin{equation*}
 H^{(1)}=H_D+[A_1,H_D]
 +\frac12[A_1,[A_1,H_D]]
 +\frac1{3!}[A_1,[A_1,[A_1,H_D]]]+\cdots .
\end{equation*}
对静质量与奇算符的嵌套对易子逐阶计算，
\begin{align*}
 [A_1,\beta_Dm_{\mathrm e} c^2]&=-\hat{\mathcal O},\\
 [A_1,\hat{\mathcal O}]&=\frac{\beta_D\hat{\mathcal O}^2}{m_{\mathrm e} c^2},\\
 [A_1,[A_1,\beta_Dm_{\mathrm e} c^2]]&=-\frac{\beta_D\hat{\mathcal O}^2}{m_{\mathrm e} c^2},\\
 [A_1,[A_1,\hat{\mathcal O}]]&=-\frac{\hat{\mathcal O}^3}{m_{\mathrm e}^2c^4}.
\end{align*}
继续到四阶并按 BCH 系数收集，原来的 $+\hat{\mathcal O}$ 被一阶项抵消，偶部分分别留下
\begin{equation*}
 \frac{\beta_D\hat{\mathcal O}^2}{2m_{\mathrm e} c^2},
 \qquad
 -\frac{\beta_D\hat{\mathcal O}^4}{8m_{\mathrm e}^3c^6}.
\end{equation*}
偶势满足 $[\beta_D,\hat{\mathcal V}]=0$。第一层对易子为奇算符，第二层对易子为
\begin{equation*}
 [A_1,[A_1,\hat{\mathcal V}]]
 =-\frac{1}{4m_{\mathrm e}^2c^4}[\hat{\mathcal O},[\hat{\mathcal O},\hat{\mathcal V}]],
\end{equation*}
乘上 BCH 的 $1/2!$ 后得到
\begin{equation*}
 -\frac{1}{8m_{\mathrm e}^2c^4}[\hat{\mathcal O},[\hat{\mathcal O},\hat{\mathcal V}]].
\end{equation*}
第一次变换仍留下三阶奇算符
\begin{equation*}
 \hat{\mathcal O}_1
 =\frac{\beta_D[\hat{\mathcal O},\hat{\mathcal V}]}{2m_{\mathrm e} c^2}
 -\frac{\hat{\mathcal O}^3}{3m_{\mathrm e}^2c^4}+\cdots .
\end{equation*}
再取 $A_2=\beta_D\hat{\mathcal O}_1/(2m_{\mathrm e} c^2)$，其与静质量的对易子正好给 $-\hat{\mathcal O}_1$，因此这一奇块被下一步变换消去。按 $\|\hat{\mathcal O}\|/(m_{\mathrm e} c^2)=O(\epsilon)$、$\|\hat{\mathcal V}\|/(m_{\mathrm e} c^2)=O(\epsilon^2)$ 计数，第二步对偶部分的首个新修正高于当前保留阶数。于是得到正文\eqnrefs{eq:fw-even-operator-form-r16}。
\end{proof}

\begin{lemma}[Dirac 奇部的平方]
若 $\mathcal O=c\boldsymbol\alpha\cdot\vb{\pi}$，则 $\mathcal O^2=c^2(\vb{\pi}^2-q\hbar\boldsymbol\Sigma\cdot\vb{B})$，其中 $\boldsymbol\Sigma$ 是本章四分量自旋矩阵。
\end{lemma}

\begin{proof}
正文\eqnrefs{eq:pi-commutator} 已说明动力学动量的反对称部分由磁场控制。对最小耦合奇算符 $\hat{\mathcal O}=c\boldsymbol\alpha\cdot\vb{\Pi}$ 求平方时，Pauli 代数的对称部分给普通动能，反对称部分恰好读取这个动量对易子，因此磁矩项由最小耦合动量的非对易部分与 Dirac--Pauli 代数共同生成。

利用
\begin{equation*}
 \alpha_i\alpha_j
 =\delta_{ij}\id()+\ii\epsilon_{ijk}\Sigma_k,
\end{equation*}
有
\begin{align*}
 \hat{\mathcal O}^2
 &=c^2\alpha_i\alpha_j\Pi_i\Pi_j,\\
 &=c^2\left[
 \vb{\Pi}^2
 +\frac{\ii}{2}\epsilon_{ijk}\Sigma_k[\Pi_i,\Pi_j]
 \right].
\end{align*}
代入正文\eqnrefs{eq:pi-commutator}，
\begin{equation*}
 \frac{\ii}{2}\epsilon_{ijk}\Sigma_k[\Pi_i,\Pi_j]
 =-q_{\mathrm e}\hbar\boldsymbol\Sigma\cdot\vb{B},
\end{equation*}
从而得到正文\eqnrefs{eq:fw-O2-r16}。
\end{proof}

\begin{lemma}[静电势的双重交换子]
在静态光滑电磁势下，$[\mathcal O,[\mathcal O,q\phi]]$ 分解为势的 Laplace 算子项与自旋轨道项，精确微分表达式为\eqnrefs{eq:fw-double-commutator-r16}。反对易对称化保证算符在本章实势条件下形式自伴。
\end{lemma}

\begin{proof}
正文\eqnrefs{eq:dirac-even-odd-decomposition-r16} 中 $\hat{\mathcal V}=q_{\mathrm e}\phi$、$\hat{\mathcal O}=c\boldsymbol\alpha\cdot\vb{\Pi}$。对静态标势，第一层对易子只依赖 $[\Pi_i,\phi]$；再与 $\hat{\mathcal O}$ 对易一次后，Pauli 代数的对称部分产生 $\nabla\cdot\vb{E}$，反对称部分产生自旋与 $\vb{E}\times\vb{\Pi}$ 的耦合。

因为 $\vb{E}=-\vb{\nabla}\phi$，
\begin{align*}
 [\hat{\mathcal O},\hat{\mathcal V}]
 &=q_{\mathrm e} c\alpha_i[\Pi_i,\phi],\\
 &=\ii q_{\mathrm e}\hbar c\,\boldsymbol\alpha\cdot\vb{E}.
\end{align*}
于是
\begin{align*}
 [\hat{\mathcal O},[\hat{\mathcal O},\hat{\mathcal V}]]
 &=\ii q_{\mathrm e}\hbar c^2[\alpha_i\Pi_i,\alpha_jE_j],\\
 &=\ii q_{\mathrm e}\hbar c^2
 \left
 \{[\Pi_i,E_i]
 +\ii\epsilon_{ijk}\Sigma_k(\Pi_iE_j+E_j\Pi_i)
 \right\}.
\end{align*}
再用 $[\Pi_i,E_i]=-\ii\hbar\partial_iE_i$ 并整理 Levi--Civita 收缩，得到
\begin{equation*}
 [\hat{\mathcal O},[\hat{\mathcal O},\hat{\mathcal V}]]
 =q_{\mathrm e}\hbar^2c^2\vb{\nabla}\cdot\vb{E}
 +q_{\mathrm e}\hbar c^2\boldsymbol\Sigma\cdot
 \left(\vb{E}\times\vb{\Pi}-\vb{\Pi}\times\vb{E}\right),
\end{equation*}
即正文\eqnrefs{eq:fw-double-commutator-r16}。
\end{proof}

\begin{proposition}[Dirac 模型的 Pauli 极限]
在低动量、静态缓变电磁背景及本节 Foldy--Wouthuysen 截断条件下，正能块的有效哈密顿量为\eqnrefs{eq:fw-pauli}，同时包含动力学能、Zeeman、自旋轨道和 Darwin 项。它是带明确截断阶的模型结论。
\end{proposition}

\begin{proof}
令
\[
P_+\equiv\frac{\id()+\beta_D}{2}.
\]
因为 $P_+\beta_D P_+=P_+$，而 $\boldsymbol\Sigma$ 在正、负能静止块内分别退化成同一组 Pauli 矩阵，所以
\[
P_+\boldsymbol\Sigma P_+
=\boldsymbol\sigma_{\rm P}P_+.
\]
对任意只作用在正能块内的态，正文\eqnrefs{eq:fw-O2-r16} 因而给出
\[
P_+\hat{\mathcal O}^2P_+
=c^2\Bigl(
\hat{\vb{\Pi}}^2
-q_{\mathrm e}\hbar\boldsymbol\sigma_{\rm P}\cdot\vb{B}
\Bigr)P_+.
\]
四次奇算符不能再单独把磁场项丢掉。由 $\hat{\mathcal O}^4=(\hat{\mathcal O}^2)^2$，有严格的算符恒等式
\begin{align*}
P_+\hat{\mathcal O}^4P_+
={}&c^4
\Bigl(
\hat{\vb{\Pi}}^2
-q_{\mathrm e}\hbar\boldsymbol\sigma_{\rm P}\cdot\vb{B}
\Bigr)^2P_+\\
={}&c^4\Bigl[
\hat{\vb{\Pi}}^4
-q_{\mathrm e}\hbar\bigl\{
\hat{\vb{\Pi}}^2,
\boldsymbol\sigma_{\rm P}\cdot\vb{B}
\bigr\}
+q_{\mathrm e}^2\hbar^2
(\boldsymbol\sigma_{\rm P}\cdot\vb{B})^2
\Bigr]P_+,
\end{align*}
其中反对易子保留了 $\hat{\vb{\Pi}}$ 与非均匀磁场不对易的信息。把这两式以及正文\eqnrefs{eq:fw-double-commutator-r16} 代入\eqnrefs{eq:fw-even-operator-form-r16}，得到
\begin{align*}
P_+H_{\rm FW}P_+
={}&\Biggl[
m_{\mathrm e} c^2+q_{\mathrm e}\phi
+\frac{\hat{\vb{\Pi}}^2-q_{\mathrm e}\hbar\boldsymbol\sigma_{\rm P}\cdot\vb{B}}{2m_{\mathrm e}}\\
&-\frac{\left(\hat{\vb{\Pi}}^2-q_{\mathrm e}\hbar\boldsymbol\sigma_{\rm P}\cdot\vb{B}\right)^2}{8m_{\mathrm e}^3c^2}
-\frac{q_{\mathrm e}\hbar^2}{8m_{\mathrm e}^2c^2}\vb{\nabla}\cdot\vb{E}\\
&-\frac{q_{\mathrm e}\hbar}{8m_{\mathrm e}^2c^2}
\boldsymbol\sigma_{\rm P}\cdot
\left(
\vb{E}\times\hat{\vb{\Pi}}
-\hat{\vb{\Pi}}\times\vb{E}
\right)
\Biggr]P_+ + P_+R_{\rm FW}P_+ .
\end{align*}
$R_{\rm FW}$ 收集前一推导中没有保留的更高阶 Foldy--Wouthuysen 项。去掉末尾的投影算符，就是正文\eqnrefs{eq:fw-pauli}。若还要使用更常见的纯 $\Pi^4$ 相对论动能修正，则必须另外控制
\[
\epsilon_{B/p}
\equiv
\frac{|q_{\mathrm e}|\hbar B_\star}{p_\star^2}\ll1;
\]
这时四次项中的反对易子与 $B^2$ 项才相对于 $\Pi^4$ 受控变小。这个附加条件与正负能块解耦条件\eqnrefs{eq:dirac-pauli-control-r16} 是不同的近似。
\end{proof}

\begin{proposition}[自由 Dirac 方程的质量壳]
设 $u\ne0$ 满足 $(c\boldsymbol\alpha\cdot\vb{p}+\beta mc^2)u=Eu$。Clifford 关系给出 $E^2=c^2\vb{p}^2+m^2c^4$，两支能量为 $\pm\sqrt{c^2\vb{p}^2+m^2c^4}$。
\end{proposition}

\begin{proof}
对正频旋量左乘 $\slashed p+m_{\mathrm e} c$：

\begin{align*}
 (\slashed p+m_{\mathrm e} c)(\slashed p-m_{\mathrm e} c)u_s
 &=\left(\gamma^\mu\gamma^\nu p_\mu p_\nu-m_{\mathrm e}^2c^2\right)u_s.
\end{align*}
把矩阵乘积分成对称与反对称两部分，
\begin{align*}
 \gamma^\mu\gamma^\nu p_\mu p_\nu
 &=\frac12\{\gamma^\mu,\gamma^\nu\}p_\mu p_\nu
 +\frac12[\gamma^\mu,\gamma^\nu]p_\mu p_\nu\\
 &=\eta^{\mu\nu}p_\mu p_\nu\,\id(4)
 =p^2\id(4),
\end{align*}
因为 $p_\mu p_\nu$ 对称，而 $[\gamma^\mu,\gamma^\nu]$ 反对称。于是
\begin{align*}
 (p^2-m_{\mathrm e}^2c^2)u_s&=0.
\end{align*}
非零旋量要求
\begin{align*}
 p^2&=m_{\mathrm e}^2c^2,\\
 \frac{E_{\vb{p}}^2}{c^2}-\vb{p}^2&=m_{\mathrm e}^2c^2,\\
 E_{\vb{p}}&=\sqrt{m_{\mathrm e}^2c^4+c^2\vb{p}^2},
\end{align*}
其中最后选择正能分支。对正文\eqnrefs{eq:v-dirac} 左乘 $\slashed p-m_{\mathrm e} c$ 得到同一二阶条件，因此电子与正电子外线共享同一正能质量壳参数化。
\end{proof}

\begin{proposition}[自由负能旋量的构造]
对 $m_{\mathrm e}>0$，固定动量的自由 Dirac 哈密顿量的负能本征空间维数为二。按本章的二分量自旋基与归一化构造旋量，可得\eqnrefs{eq:explicit-v}。本章 $v_s(p)$ 的负频参数使用正能四动量 $p$，故它属于 $H(-\vb{p})$ 的负能空间；普通厄米内积的正交关系为 $u_s(p)^\dagger v_{s'}(E_{\vb{p}}/c,-\vb{p})=0$。
\end{proposition}

\begin{proof}
在 Dirac 表示中把四分量旋量写成两个二分量块。对正频支令
\[
u_s(p)=\begin{pmatrix}\phi_s\\ \chi_s\end{pmatrix}.
\]
正文\eqnrefs{eq:u-dirac} 给出
\begin{align*}
(E-m_{\mathrm e} c^2)\phi_s
-c\boldsymbol\sigma_{\rm P}\!\cdot\!\vb{p}\,\chi_s&=0,\\
c\boldsymbol\sigma_{\rm P}\!\cdot\!\vb{p}\,\phi_s
-(E+m_{\mathrm e} c^2)\chi_s&=0.
\end{align*}
第二式先确定下分量：
\[
\chi_s
=\frac{c\boldsymbol\sigma_{\rm P}\cdot\vb{p}}{E+m_{\mathrm e} c^2}\phi_s
=\frac{\boldsymbol\sigma_{\rm P}\cdot\vb{p}}{E/c+m_{\mathrm e} c}\phi_s.
\]
取
$\phi_s=\sqrt{E/c+m_{\mathrm e} c}\,\xi_s^{(+)}$，便得到正文\eqnrefs{eq:explicit-u}。用质量壳 $E^2=m_{\mathrm e}^2c^4+c^2\vb{p}^2$ 检查归一化：
\begin{align*}
\bar u_su_{s'}
&=\left(\frac Ec+m_{\mathrm e} c\right)
\left[
1-\frac{\vb{p}^2}{(E/c+m_{\mathrm e} c)^2}
\right]\delta_{ss'}\\
&=2m_{\mathrm e} c\,\delta_{ss'}.
\end{align*}

对负频支令
\[
v_s(p)=\begin{pmatrix}\phi'_s\\ \chi'_s\end{pmatrix}.
\]
正文\eqnrefs{eq:v-dirac} 的两个分量方程为
\begin{align*}
(E+m_{\mathrm e} c^2)\phi'_s
-c\boldsymbol\sigma_{\rm P}\!\cdot\!\vb{p}\,\chi'_s&=0,\\
c\boldsymbol\sigma_{\rm P}\!\cdot\!\vb{p}\,\phi'_s
-(E-m_{\mathrm e} c^2)\chi'_s&=0.
\end{align*}
第一式给
\[
\phi'_s
=\frac{\boldsymbol\sigma_{\rm P}\cdot\vb{p}}{E/c+m_{\mathrm e} c}\chi'_s.
\]
取
$\chi'_s=\sqrt{E/c+m_{\mathrm e} c}\,\xi_s^{(-)}$，即得到正文\eqnrefs{eq:explicit-v}。同样利用质量壳，
\begin{align*}
\bar v_sv_{s'}
&=\left(\frac Ec+m_{\mathrm e} c\right)
\left[
\frac{\vb{p}^2}{(E/c+m_{\mathrm e} c)^2}-1
\right]\delta_{ss'}\\
&=-2m_{\mathrm e} c\,\delta_{ss'}.
\end{align*}
最后将 $v$ 的空间动量反向，记 $K=\boldsymbol\sigma_{\rm P}\cdot\vb{p}/(E/c+m_{\mathrm e}c)$，则 $K^\dagger=K$。两旋量的分块正比于 $(\xi,K\xi)$ 与 $(-K\eta,\eta)$，因此
\[
u_s(p)^\dagger v_{s'}(E/c,-\vb{p})
=\left(E/c+m_{\mathrm e}c\right)(-\xi_s^\dagger K\eta_{s'}+\xi_s^\dagger K\eta_{s'})=0.
\]
\end{proof}

\begin{proposition}[自由旋量的完备关系]
设 $p^2=m_{\mathrm e}^2c^2$、$p^0>0$，采用本章 $u,v$ 的归一化，则
\begin{align*}
\Lambda_+(p)&=\frac{\sum_su_s(p)\bar u_s(p)}{2m_{\mathrm e}c}
=\frac{\slashed p+m_{\mathrm e}c}{2m_{\mathrm e}c},\\
\Lambda_-(p)&=-\frac{\sum_sv_s(p)\bar v_s(p)}{2m_{\mathrm e}c}
=\frac{m_{\mathrm e}c-\slashed p}{2m_{\mathrm e}c}.
\end{align*}
它们是 $\slashed p$ 的两个协变代数投影，满足 $\Lambda_\pm^2=\Lambda_\pm$、$\Lambda_+\Lambda_-=0$、$\Lambda_++\Lambda_-=I_4$。它们在普通厄米内积下通常是斜投影；同一动量的哈密顿量的能量正交投影须另行定义。未归一化的自旋求和公式即\eqnrefs{eq:u-complete}与\eqnrefs{eq:v-complete}。
\end{proposition}

\begin{proof}
质量壳给 $\slashed p^{\,2}=m_{\mathrm e}^2c^2$，因此
\begin{align*}
\Lambda_+^2
&=\frac{(\slashed p+m_{\mathrm e} c)^2}{4m_{\mathrm e}^2c^2}
=\frac{\slashed p+m_{\mathrm e} c}{2m_{\mathrm e} c}=\Lambda_+,\\
\Lambda_-^2
&=\frac{(m_{\mathrm e} c-\slashed p)^2}{4m_{\mathrm e}^2c^2}
=\Lambda_-.
\end{align*}
并且
\begin{equation*}
(\slashed p-m_{\mathrm e} c)\Lambda_+=0,
\qquad
(\slashed p+m_{\mathrm e} c)\Lambda_-=0.
\end{equation*}
由于 $\operatorname{tr}\slashed p=0$，两投影的迹均为 $2$，所以其像空间分别与正能、负频 Dirac 解空间完全相同。再由伴随 Dirac 方程 $\bar u_s\slashed p=m_{\mathrm e}c\bar u_s$ 及 $\slashed p v_{s'}=-m_{\mathrm e}cv_{s'}$ 得 $\bar u_sv_{s'}=0$；同理 $\bar v_su_{s'}=0$。因此归一化基求和在一个特征空间上恒等，在另一个特征空间上为零，与相应 $\Lambda_\pm$ 的作用完全相同。由 $\bar u_su_{s'}=2m_{\mathrm e} c\,\delta_{ss'}$，正能解空间上的恒等投影为
\begin{equation*}
\Lambda_+=\frac1{2m_{\mathrm e} c}\sum_su_s\bar u_s,
\end{equation*}
故得到\eqnrefs{eq:u-complete}。负频基满足 $\bar v_sv_{s'}=-2m_{\mathrm e} c\,\delta_{ss'}$，于是
\begin{equation*}
\Lambda_-=-\frac1{2m_{\mathrm e} c}\sum_sv_s\bar v_s,
\end{equation*}
等价于\eqnrefs{eq:v-complete}。
\end{proof}

\begin{proposition}[电子色散 Taylor 截断的控制]
在固定中心动量的有限邻域内，精确自由色散 $E(p)=\sqrt{m^2c^4+c^2p^2}$ 的二阶近似误差由三阶导数的上确界乘 $|\Delta p|^3/6$ 控制；传播时还须使该能量误差乘作用时间除以 $\hbar$ 足够小。该判据为\eqnrefs{eq:third-order-control}。
\end{proposition}

\begin{proof}
为得到以动量宽度 $\sigma_p$ 表示的局部截断判据，取工作窗边界 $|\delta p|=\sigma_p$；正文\eqnrefs{eq:dispersion-taylor} 的最后保留项和第一个舍弃项分别为
\[
\Delta E_2=\frac12E''(p_0)\sigma_p^2,
\qquad
\Delta E_3=\frac16E^{(3)}(p_0)\sigma_p^3.
\]
先从精确色散 $E(p)=\sqrt{m_{\mathrm e}^2c^4+p^2c^2}$ 计算所需导数：
\begin{align*}
E'(p)&=\frac{pc^2}{E(p)},\\
E''(p)&=\frac{m_{\mathrm e}^2c^6}{E(p)^3},\\
E^{(3)}(p)
&=-\frac{3m_{\mathrm e}^2c^8p}{E(p)^5}.
\end{align*}
在 $p=p_0$ 处用 $E_0=\gamma_0m_{\mathrm e} c^2$ 与 $p_0c/E_0=\beta_0$，两项之比为
\begin{align*}
\epsilon_{\rm curv}
&=\left|\frac{\Delta E_3}{\Delta E_2}\right|\\
&=\frac{|E^{(3)}(p_0)|\sigma_p}{3|E''(p_0)|}\\
&=\frac{p_0c^2}{E_0^2}\sigma_p\\
&=\frac{p_0^2c^2}{E_0^2}\frac{\sigma_p}{p_0}\\
&=\beta_0^2\frac{\sigma_p}{p_0}.
\end{align*}
这就是正文\eqnrefs{eq:third-order-control}。
\end{proof}

\begin{proposition}[质量壳上的反冲相位匹配]
若入射和出射电子满足同一质量壳，且按\eqnrefs{eq:four-transfer}定义传递四动量，则两条质量壳相减给出\eqnrefs{eq:recoil-phase-match}。该恒等式未预先忽略二次反冲项。
\end{proposition}

\begin{proof}
正文\eqnrefs{eq:four-transfer} 把正频模式与能量流向分开写成
\[
P_f=P_i+\sigma\hbar K,
\qquad \sigma^2=1.
\]
初态和末态都是同一静质量的真实电子，所以
\[
P_i^2=P_f^2=m_{\mathrm e}^2c^2.
\]
两式相减时静质量项完全消去，而 $\sigma^2=1$ 保留二阶反冲项：
\begin{align*}
0
&=P_f^2-P_i^2\\
&=(P_i+\sigma\hbar K)^2-P_i^2\\
&=2\sigma\hbar P_i\!\cdot K+\hbar^2K^2.
\end{align*}
除以 $\hbar$ 得到正文\eqnrefs{eq:exact-recoil-four}。因此吸收与发射的差别只出现在与 $P_i\!\cdot K$ 相乘的一阶流向符号中；模式本身仍始终取 $\omega>0$。

在一维纵向通道中，$P_f=P_i+\sigma\hbar K$ 给出
\begin{align*}
p_f&=p_0+\sigma\hbar\kappa_z,\\
E_f-E_i&=\sigma\hbar\omega,
\end{align*}
故正文\eqnrefs{eq:exact-longitudinal-match} 是四维母恒等式的能量表示。若 $|\hbar\kappa_z|$ 落在正文已经给出的 Taylor 控制窗内，对
$\delta p=\sigma\hbar\kappa_z$ 使用带积分余项的 Taylor 定理：
\begin{align*}
E(p_0+\delta p)-E(p_0)
&=E'(p_0)\delta p
+\frac{E''(p_0)}{2}\delta p^2
+\frac{\delta p^3}{2}
\int_0^1\dd s\,(1-s)^2E^{(3)}(p_0+s\delta p)\\
&=\sigma v_0\hbar\kappa_z
+\frac{\hbar^2\kappa_z^2}{2m_\parallel}
+\mathcal R_3^{(\sigma)}(\kappa_z).
\end{align*}
与 $\sigma\hbar\omega$ 比较，并利用 $1/\sigma=\sigma$，得到正文\eqnrefs{eq:recoil-phase-match}。积分余项还立刻给出
\begin{align*}
\left|\mathcal R_3^{(\sigma)}(\kappa_z)\right|
&\le
\frac{|\hbar\kappa_z|^3}{2}
\sup_{|p-p_0|\le\hbar|\kappa_z|}|E^{(3)}(p)|
\int_0^1(1-s)^2\dd s\\
&=\frac{\hbar^3|\kappa_z|^3}{6}
\sup_{|p-p_0|\le\hbar|\kappa_z|}|E^{(3)}(p)|.
\end{align*}
这里可以直接看出：线性同步项不区分吸收和发射，第一项真正区分两种流向的反冲修正出现在偶次色散项；舍弃高阶项时也有明确的误差上界，而不是只靠一个形式上的省略号。
\end{proof}

\begin{proposition}[Maxwell 散射的 T 算符]
若背景 Green 算符 $G_0$ 已固定，扰动算符为 $V$ 且 $I-G_0V$ 可逆，则 Lippmann--Schwinger 方程有唯一解，散射部分可写为 $G_0T$ 作用于入射场，其中 $T=V(I-G_0V)^{-1}$。表达式不要求 Born 展开收敛。
\end{proposition}

\begin{proof}
把背景 Maxwell Green 算符记为 $\mathbf G_0^R$，介质相对于背景的散射势记为 $\mathsf V_{\rm EM}$。正文的 Lippmann--Schwinger 方程是
\[
\widetilde{\vb{E}}
=\widetilde{\vb{E}}_{\rm inc}
+\mathbf G_0^R\mathsf V_{\rm EM}\widetilde{\vb{E}}.
\]
先不要展开级数，而把未知总场的项移到左边：
\[
(\mathbf I-\mathbf G_0^R\mathsf V_{\rm EM})\widetilde{\vb{E}}
=\widetilde{\vb{E}}_{\rm inc}.
\]
只要所考察复频率不位于该算符的谱上，就存在逆算符，故
\[
\widetilde{\vb{E}}
=(\mathbf I-\mathbf G_0^R\mathsf V_{\rm EM})^{-1}
\widetilde{\vb{E}}_{\rm inc}.
\]
散射场因此为
\begin{align*}
\widetilde{\vb{E}}_{\rm sc}
&\equiv\widetilde{\vb{E}}-\widetilde{\vb{E}}_{\rm inc}\\
&=(\mathbf I-\mathbf G_0^R\mathsf V_{\rm EM})^{-1}
\mathbf G_0^R\mathsf V_{\rm EM}\widetilde{\vb{E}}_{\rm inc}.
\end{align*}
现在用可直接验证的算符恒等式
\[
(\mathbf I-AB)^{-1}A=A(\mathbf I-BA)^{-1}.
\]
令 $A=\mathbf G_0^R$、$B=\mathsf V_{\rm EM}$，便有
\[
\widetilde{\vb{E}}_{\rm sc}
=\mathbf G_0^R
\mathsf V_{\rm EM}
(\mathbf I-\mathbf G_0^R\mathsf V_{\rm EM})^{-1}
\widetilde{\vb{E}}_{\rm inc}.
\]
因此定义
\[
{
\mathsf T_{\rm EM}
\equiv\mathsf V_{\rm EM}
(\mathbf I-\mathbf G_0^R\mathsf V_{\rm EM})^{-1}}
\]
后立即得到
\[
{
\widetilde{\vb{E}}_{\rm sc}
=\mathbf G_0^R\mathsf T_{\rm EM}\widetilde{\vb{E}}_{\rm inc}},
\]
即正文\eqnrefs{eq:maxwell_scattered_field_t_operator}。

多次散射级数只是这个精确逆算符在收敛区域内的 Neumann 展开。若某个算符范数满足 $\lVert\mathbf G_0^R\mathsf V_{\rm EM}\rVert<1$，则
\begin{align*}
(\mathbf I-\mathbf G_0^R\mathsf V_{\rm EM})^{-1}
&=\sum_{n=0}^{\infty}(\mathbf G_0^R\mathsf V_{\rm EM})^n,\\
\mathsf T_{\rm EM}
&=\mathsf V_{\rm EM}
+\mathsf V_{\rm EM}\mathbf G_0^R\mathsf V_{\rm EM}
+\mathsf V_{\rm EM}\mathbf G_0^R\mathsf V_{\rm EM}\mathbf G_0^R\mathsf V_{\rm EM}+\cdots.
\end{align*}
第 $n$ 个额外因子 $\mathbf G_0^R\mathsf V_{\rm EM}$ 对应“再传播一次、再散射一次”。即使 Neumann 级数本身不收敛，$\mathsf T_{\rm EM}$ 仍由前面的算符逆定义；其极点对应齐次方程
\[
(\mathbf I-\mathbf G_0^R\mathsf V_{\rm EM})\vb{E}=0
\]
出现非零解，也就是
$\ker(\mathbf I-\mathbf G_0^R\mathsf V_{\rm EM})\neq\{0\}$；在有限维离散化后，这才等价于相应行列式为零。它对应无外源条件下仍可存在的本征模或共振。这样 $T$ 算符同时统一了弱散射 Born 级数与强多重散射的谱结构。
\end{proof}

\begin{theorem}[互易介质的 Green 对称性]
对满足本章双线性互易条件的线性介质，在相同边界条件及辐射条件下，电磁 Green 张量满足\eqnrefs{eq:maxwell-green-reciprocity-r16}的源观察点交换与转置关系。一般磁光非互易介质不满足此前提。
\end{theorem}

\begin{proof}
正文\eqnrefs{eq:maxwell-bilinear-reciprocity-dp68} 已经把介质互易性写成 Maxwell 算符的双线性对称关系。要把它变成 Green 张量的源点交换关系，只需让两组测试场分别由两个点源产生。

取任意常矢量 \(\vb{a},\vb{b}\)，令
\begin{equation*}
\vb{u}(\vb{s})=\mathbf G^R(\vb{s},\vb{r};\omega)\cdot\vb{a},
\qquad
\vb{v}(\vb{s})=\mathbf G^R(\vb{s},\vb{r}';\omega)\cdot\vb{b}.
\end{equation*}
由 Green 方程，
\begin{equation*}
\mathsf L_{\rm EM}\vb{u}
=\vb{a}\,\delta^{(3)}(\vb{s}-\vb{r}),
\qquad
\mathsf L_{\rm EM}\vb{v}
=\vb{b}\,\delta^{(3)}(\vb{s}-\vb{r}').
\end{equation*}
代入正文\eqnrefs{eq:maxwell-bilinear-reciprocity-dp68} 后，两个空间积分分别被 \(\delta\) 函数取样：
\begin{equation*}
\vb{a}\cdot\mathbf G^R(\vb{r},\vb{r}';\omega)\cdot\vb{b}
=\vb{b}\cdot\mathbf G^R(\vb{r}',\vb{r};\omega)\cdot\vb{a}.
\end{equation*}
因为 \(\vb{a},\vb{b}\) 任意，所以
\begin{equation*}
\mathbf G^R(\vb{r},\vb{r}';\omega)
=[\mathbf G^R(\vb{r}',\vb{r};\omega)]^T,
\end{equation*}
即正文\eqnrefs{eq:maxwell-green-reciprocity-r16}。
\end{proof}

\begin{proposition}[电磁 Green 算符的光学恒等式]
对本章因果 Maxwell 算符，在共同定义域及一致辐射边界条件下，$G-G^\dagger=G(L^\dagger-L)G^\dagger$。其含耗散介质和边界通道的电磁表达为\eqnrefs{eq:green-optical-identity-SI-r12}。
\end{proposition}

\begin{proof}
先在整个外部空间加入正的被动吸收。对固定实频 $\omega>0$，场在无穷远衰减，因此旋度算符的分部积分表面项消失，$\mathsf L_{\rm EM}$ 与其伴随可按同一完整空间处理。以下代数在该正则化下完成；弱吸收极限必须放在空间积分之后。定义
\[
\mathbf G^R=\mathsf L_{\rm EM}^{-1},
\qquad
\mathbf G^{R\dagger}=(\mathsf L_{\rm EM}^{\dagger})^{-1}.
\]
对任意两个可逆算符 $A,B$，有两种等价的预解式恒等式
\[
A^{-1}-B^{-1}
=A^{-1}(B-A)B^{-1}
=B^{-1}(B-A)A^{-1}.
\]
第一式例如由
$A^{-1}BB^{-1}-A^{-1}AB^{-1}=A^{-1}-B^{-1}$ 逐项得到。取
$A=\mathsf L_{\rm EM}$、$B=\mathsf L_{\rm EM}^{\dagger}$，先用第一种次序得到
\begin{align*}
\mathbf G^R-\mathbf G^{R\dagger}
&=\mathbf G^R
(\mathsf L_{\rm EM}^{\dagger}-\mathsf L_{\rm EM})
\mathbf G^{R\dagger}.
\end{align*}
在上述外部完成下，材料和外部吸收信息集中在算符差 $\mathsf L_{\rm EM}^{\dagger}-\mathsf L_{\rm EM}$ 中。对局域电介质且磁响应无损的情形，
\[
\mathsf L_{\rm EM}
=\nabla\times\nabla\times
-\frac{\omega^2}{c^2}\boldsymbol\varepsilon_r,
\]
故
\[
\mathsf L_{\rm EM}^{\dagger}-\mathsf L_{\rm EM}
=\frac{\omega^2}{c^2}
(\boldsymbol\varepsilon_r-\boldsymbol\varepsilon_r^{\dagger})
=2\ii\frac{\omega^2}{c^2}
\operatorname{Im}\boldsymbol\varepsilon_r.
\]
代回并除以 $2\ii$：
\[
\operatorname{Im}\mathbf G^R
=\frac{\omega^2}{c^2}
\mathbf G^R
(\operatorname{Im}\boldsymbol\varepsilon_r)
\mathbf G^{R\dagger}.
\]
把算符乘积写成坐标核。若介电张量局域，
$\operatorname{Im}\boldsymbol\varepsilon_r$ 的核为
$\operatorname{Im}\boldsymbol\varepsilon_r(\vb{s},\omega)
\delta^{(3)}(\vb{s}-\vb{s}')$，于是
\[
{
\operatorname{Im}\mathbf G^R(\vb{r},\vb{r}';\omega)
=\frac{\omega^2}{c^2}\int\dd^3s\,
\mathbf G^R(\vb{r},\vb{s};\omega)
\cdot\operatorname{Im}\boldsymbol\varepsilon_r(\vb{s},\omega)
\cdot[\mathbf G^R(\vb{r}',\vb{s};\omega)]^\dagger}.
\]
这就是正文\eqnrefs{eq:green-optical-identity-SI-r12}。

该恒等式还解释了为什么 $\operatorname{Im}\mathbf G^R$ 能作为损失谱权重。为把正定性写成最直观的场能量形式，改用上面第二种等价次序，亦有
\[
\operatorname{Im}\mathbf G^R
=\frac{\omega^2}{c^2}
\mathbf G^{R\dagger}
(\operatorname{Im}\boldsymbol\varepsilon_r)
\mathbf G^R.
\]
对任意测试电流 $\vb{J}$，定义由它激发的场
\[
\vb{E}_J(\vb{s})
\equiv\int\dd^3r'\,\mathbf G^R(\vb{s},\vb{r}';\omega)\cdot\vb{J}(\vb{r}').
\]
于是左右与 $\vb{J}^*,\vb{J}$ 收缩得到
\[
\iint\dd^3r\dd^3r'\,
\vb{J}^*(\vb{r})\cdot\operatorname{Im}\mathbf G^R(\vb{r},\vb{r}')\cdot\vb{J}(\vb{r}')
=\frac{\omega^2}{c^2}\int\dd^3s\,
\vb{E}_J^*(\vb{s})\cdot
\operatorname{Im}\boldsymbol\varepsilon_r(\vb{s})\cdot\vb{E}_J(\vb{s}).
\]
对被动介质，$\operatorname{Im}\boldsymbol\varepsilon_r\succeq0$，所以右边非负。因而 Green 张量的反厄米部分不是一个形式上的“取虚部”规则，而是由材料局域耗散经过两次传播映射回源空间后得到的正半定谱核；这正是后文 EELS、LDOS 与涨落--耗散关系共同使用的数学结构。


若外部保持严格真空，入射与出射定义域不相同，应另计边界通量；该通量正是正文\eqnrefs{eq:feqo-radiative-spectral-operator} 的 $\mathsf R_{\rm rad}$。真空点偶极的 $\operatorname{Im}G_0=kI/(6\pi)$ 说明该项即使没有任何体吸收也不会消失。
\end{proof}
